DGFiP, IRCOM et déclarations 2042, revenus 2024 ; CCSS ; calculs de l'auteur.\par Note : 98 barèmes « franchise (8 à 20 k\euro{}), taux unique jusqu'à 60-200 k\euro{}, 40 ou 45 \%, 49 \% au-delà de 250 k\euro{} », tous calés sur la même recette ; scénario central. Ordonnée : plus faible écart de Lorenz aux bornes des tranches (positif : dominance). Les points bleus remplissent aussi la condition de hausse limitée à 2 points ; écarts inférieurs à $-$12 ‰ ramenés à $-$12.
